\documentclass{article}

\usepackage[T2A]{fontenc}
\usepackage[utf8]{inputenc}
\usepackage[english,russian]{babel}
\usepackage[rus]{datetime2}          \usepackage{etoolbox}                \usepackage{amsmath, amssymb}
\usepackage{enumitem}
\usepackage{tikz}
\usepackage{graphicx}
\usepackage{stackengine}
\usepackage{hyperref}
\hypersetup{
    linktocpage,
    colorlinks=true
}

\newcounter{theorem}
\newcounter{lemma}[theorem]   \newcounter{definition}

\newenvironment{theorem}[1][]{\par\medskip\noindent
  \refstepcounter{theorem}\textbf{Теорема \thetheorem.}\ifx\relax#1\relax\else\ {\normalfont(#1)}\fi
  \par\nobreak\noindent\ignorespaces
}{\par\medskip
}

\newenvironment{lemma}[1][]{\par\medskip\noindent
  \refstepcounter{lemma}\textbf{Лемма \thelemma.}\ifx\relax#1\relax\else\ {\normalfont(#1)}\fi
  \par\nobreak\noindent\ignorespaces
}{\par\medskip
}

\newenvironment{definition}[1][]{\par\medskip\noindent
  \refstepcounter{definition}\textbf{Определение \thedefinition.}\ifx\relax#1\relax\else\ {\normalfont(#1)}\fi
  \par\nobreak\noindent\ignorespaces
}{\par\medskip
}

\newenvironment{proof}{\par\medskip\noindent
  \textbf{Доказательство. }\par\nobreak\noindent\ignorespaces
}{\par\nobreak
  \noindent\hfill$\square$\par\medskip
}

\newenvironment{note}{\par\medskip\noindent
  \textbf{Замечание. }\par\nobreak\noindent\ignorespaces
}{\par\medskip
}

\title{Уравнения математической физики\\Лекционный материал}
\author{}
\date{Скомпилирован: \today}

\begin{document}

\maketitle

\newpage

\section{Лекция 1: Уравнения в частных производных 2-го порядка}

\subsection{Примеры}

1) распростронение волн - волновое уравнение \par\medskip
\begin{tabular}{p{0.45\textwidth} p{0.45\textwidth}}
  \hspace{1cm}\(\displaystyle 
    \frac{\partial^2 u}{\partial t^2} = c^2 \Delta u + \underbrace{f(x, y, z, t)}_{\text{источник}},
  \) &
  \(c\) — скорость распространения
\end{tabular}  \bigskip\\
2) теплопередача - уравнение теплопроводности \par\medskip
\begin{tabular}{p{0.45\textwidth} p{0.45\textwidth}}
  \hspace{1cm}\(\displaystyle 
    \frac{\partial u}{\partial t} = c^2 \Delta u + \underbrace{f(x, y, z, t)}_{\text{источник}} ,
  \) &
  $c=\frac{k}{\rho C}$ -- теплопроводность
\end{tabular} \bigskip\\
3) уравнение Пуассона-Лапласа:\par\medskip
\begin{tabular}{p{0.45\textwidth} p{0.45\textwidth}}
  \hspace{1cm}\(\displaystyle 
    \Delta u = -\tilde f(x, y, z, t) ,
  \) &
  $\tilde f=\frac{f}{a^2}$ -- приведенное уравнение
\end{tabular} \\
\begin{tabular}{p{0.45\textwidth} p{0.45\textwidth}}
  \hspace{1cm}\(\displaystyle 
    \Delta u = 0
  \) &
\end{tabular} \\

Это основные уравнения математической физики. Они имеют бесчисленное множество решений.

\subsection{Постановка задачи}
Дополнительные условия.
\begin{enumerate}[nosep]
  \item граничные условия(краевые)
  \item начальныые условия
\end{enumerate}

\begin{definition}[Корректность постановки задач математической физики(МФ)]
Задача МФ называется поставленной корретно (по Адамару), если её решение:
\begin{enumerate}[nosep]
  \item существует;
  \item единственно;
  \item устойчиво (непрерывно зависит от данных параметров) - условие физической определенности задачи.
\end{enumerate}
\end{definition}

\label{abbr:LDE2O}
\begin{definition}[Уравнение в частных производных второго порядка(УЧП2П):]
Пусть $G$ --- область в $\mathbb{R}^n$, x=($x_1$, \ldots, $x_n$) --- точка в $\mathbb{R}^n$. УЧП2П называется соотношение, связывающее неизвестные функции n переменных u=u(x), и её производных 1\textsuperscript{го} и 2\textsuperscript{го} порядка:
\[
\frac{\partial u}{\partial x_1}, \ldots, \frac{\partial u}{\partial x_n}, \frac{\partial^2 u}{\partial^2 x_1}, \ldots, \frac{\partial^2 u}{\partial^2 x_n}
\] вида:
\[
F(x_1, \ldots x_n, \frac{\partial u}{\partial x_1}, \ldots, \frac{\partial u}{\partial x_n}, \frac{\partial^2 u}{\partial^2 x_1}, \ldots, \frac{\partial^2 u}{\partial^2 x_n}, u)=0
\]
\end{definition}
\newcommand{\LDUSP}{\hyperref[abbr:LDE2O]{ЛДУ2П}}


Решением \LDUSP{} нзывается функция класса $u(x) \in C^2(G): u(x)$ удовлетворяет данному уравнению.

\begin{note}
Замечание. здесь мы рассматриваем линейные уравнения относительно старших производных (УЧП2ПЛОСП)
\[
\sum_{i,j=1}^{n} a_{ij}(x)\frac{\partial^2u}{\partial x_i \partial x_j} + \Phi(x_1, \ldots, x_n, u, \operatorname{grad}(u)) = 0
\]
Где $a_{ij}(x) \in C(G)$ -- старшие коэффициенты $i, j = \overline{1,n}$ и\\
\hspace*{\fill}  $\Phi(x_1, \ldots, x_n, u, \operatorname{grad}(u))$ -- некоторая функция.
\end{note}

\begin{note}
$\operatorname{grad}(u) = (\frac{\partial u}{\partial x_1}, \ldots, \frac{\partial u}{\partial x_n})$.
\end{note}
\begin{note}
Если $\Phi$ - линейная относительно $u(x)$, $\operatorname{grad}(u)$, то получаем линейное уравнение в частных производных 2\textsuperscript{го} порялка(ЛУЧП2П):
$$\sum_{i,j=1}^{n} a_{ij}(x)\frac{\partial^2u}{\partial x_i \partial x_j} + \sum_{i=1}^{n} b_{i}(x)\frac{\partial u}{\partial x_i} + c(x)u = f(x),$$
\[
  x \in G,\quad a_{ij}(x), b_i(x), c(x) \in C(G); \qquad\qquad (i,j = \overline{1,n})
\]
$$f(x) \in C(G)$$
\end{note}

Иногда пишут:
$L[u]=f(x)$, где
$$\displaystyle L= \sum_{i,j=1}^{n} a_{ij}(x)\frac{\partial^2}{\partial x_i \partial x_j} + \sum_{i=1}^{n} b_{i}(x)\frac{\partial }{\partial x_i} + c(x)$$
$\displaystyle L: C^2(G)\mapsto C(G)$ -- линейный дифференциальный оператор 2\textsuperscript{го} порядка(ЛДО2П).\par\medskip


Если же все коэффициенты ЛДУЧП2П: не зависят от $\vec x$, то получаем ЛДУЧП2П с постоянными коэффициентами(ЛДУЧП2ППК):

\[
\sum_{i,j=1}^{n} a_{ij}\frac{\partial^2u}{\partial x_i \partial x_j} + \sum_{i=1}^{n} b_{i}\frac{\partial u}{\partial x_i} + cu = f(x), \quad x \in G,
\]
где \(a_{ij}, b_i, c\) --- коэффициенты, \(i,j = \overline{1,n}\), \(f(x) \in C(G)\) --- правая часть. Это уравнение можно записать как \(L[u]=f\).

Здесь
$\displaystyle L= 
  \underbrace{\sum_{i,j=1}^{n} a_{ij}\frac{\partial}{\partial x_i} \frac{\partial}{\partial x_j}}_{\text{квадратичная форма}}
  +\underbrace{\sum_{i=1}^{n} b_{i}\frac{\partial }{\partial x_i}}_{\text{линейная форма}} + c$

\subsection{Линейная невырожденная замена переменных}
Осуществим в ЛДУЧП2ППК линейную невырожденную замену переменных:

$\displaystyle\vec x=\vec x(\vec \xi),$ где $\displaystyle\vec \xi=(\xi_1, \ldots, \xi_n),\quad \xi_k = \sum_{i=1}^{n}s_{ik}x_i,\ (k=\overline{1,n})$.

В матричной форме: $$\vec \xi = S^T \vec x$$

S --- матрица коэффициентов: S = ($s_{ik}$)$\Rightarrow$

\[
\left.
\begin{array}{rcl}
\xi_1 &=& s_{11}x_1 + \ldots + s_{n1}x_n \\
      & & \ldots \\
\xi_i &=& s_{1i}x_1 + \ldots + s_{ni}x_n \\
      & & \ldots \\
\xi_n &=& s_{1n}x_1 + \ldots + s_{nn}x_n
\end{array}
\right\}
\det S \neq 0
\quad\Longrightarrow\quad
\text{$\vec x=(S^{-1})^T\vec\xi$}
\]

Обозначим
\[
u(x(\xi))=\tilde{u}(\xi) \Rightarrow \tilde{u}(\xi)=u(x)\]
\[
\text{Т.к. } \frac{\partial \xi_k}{\partial x_i} = s_{ik} \text{, то} \frac{\partial u}{\partial x_i} = \sum_{k=1}^n\frac{\partial \tilde{u}}{\partial \xi_i}\frac{\partial \xi_k}{\partial x_i} = \sum_{k=1}^ns_{ik}\frac{\partial \tilde{u}}{\partial \xi_k}
\]


\begin{note}
Из $\displaystyle\frac{\partial }{\partial x_i} = \sum_{k=1}s_{ik}\frac{\partial}{\partial \xi_k}$ имеем:
\[\frac{\partial^2u}{\partial x_i \partial x_j}=\frac{\partial}{\partial x_i}(\frac{\partial u}{\partial x_j})=\sum_{l,k=1}^{n}\frac{\partial^2 \tilde{u}}{\partial\xi_k \partial\xi_l}s_{ik}s_{jl}\]
Подставляя полученные выражения для $\frac{\partial u}{\partial x_i}$ и $\frac{\partial^2 u}{\partial x_i \partial x_j}$ в ЛДУЧП2ППК, получим:

\[\xi_k=\sum_{i=1}^{n}s_{ik}x_i,\ \xi=S^Tx, x=(S^{-1})^T\xi\Rightarrow\]

\[
\sum_{k,l=1}^n \frac{\partial^2 u}{\partial \xi_i \partial \xi_j} \sum_{i,j=1}^n a_{ij}s_{ik}s_{jl} + \sum_{k=1}^n \frac{\partial \tilde{u}}{\partial \xi_k}\sum_{i=1}^nb_is_{ik}+c\tilde{u}=\tilde{f}
\]

Обозначим через 
\[\tilde{a}_{kl}=\sum_{i,j=1}^na_{ij}s_{ik}s_{jl}, \quad\tilde{b}_k=\sum_{i=1}^nb_is_{ik}\]
получаем:
$\displaystyle
\tilde{L}\tilde{u}=\sum_{k,l=1}^n\tilde{a}_{kl}\frac{\partial^2\tilde{u}}{\partial\xi_k\partial\xi_l}+\sum_{k=1}^n\tilde{b}_k\frac{\partial\tilde{u}}{\xi_k}+c\tilde{u}=\tilde{f}(\xi),
$

где $\tilde{f}(\xi)=f(x(\xi))$

При этом матрица $A=(a_{ik})$ старших коэффициентов исходного ЛУЧП2ППК преобразуется по формулам\\
$
\displaystyle
  \tilde{A}=S^TAS 
$
\quad (из линейной алгебры)\\
Т.е. при невырожденной линейной замене переменных:
\[
\xi=S^Tx, detS\ne0,
\] матрица A преобразуется по формулам
\[
K=\sum_{i,j=1}^na_{ij}y_iy_j=y^TAy,\ y=(y_1, \ldots, y_n)^T
\]
\end{note}

\begin{note}
считаем A --- симметричной ($A^T=A$)
\end{note}
 
\end{document}
